% RESULTADO_RECUPERADO desde II REV09. Véase MAPA_PROCEDENCIA.json.
% Selección de líneas y cambios mecánicos explícitos; ninguna prueba nueva.
\subsection{Normalisation de l'état et signification du Hamiltonien modulaire}
La deuxième observable est explicitée par l'état qui la définit.
La décomposition tensorielle précédente est maintenue : dix-huit qutrits
par secteur, chacun muni de l'opérateur nombre $\operatorname{diag}(0,1,2)$ ;
$N_m$ est la somme de ces opérateurs et les secteurs agissent sur des facteurs
distincts. Soit
\[
 z_m=1+e^{-m\Delta_{\rm mod}}+e^{-2m\Delta_{\rm mod}},
 \qquad m\in\{90,120\}.
\]
La famille d’états de la source a la réalisation produit
\[
 \rho_A=Z^{-1}e^{-\Delta_{\rm mod}(90N_{90}+120N_{120})},
 \qquad Z=z_{90}^{18}z_{120}^{18}.
\]
La commutativité des opérateurs nombre et la factorisation tensorielle
prouvent l'expression de $Z$. L'état est de rang plein pour
$\Delta_{\rm mod}$ réel et fini, et son logarithme spectral donne
\[
 K_A=-\log\rho_A=(\log Z)I+\Delta_{\rm mod}D.
\]
Pour $\Delta_{\rm mod}\ne0$, $K_A$ pondère différemment les excitations
des deux secteurs. La lecture conjointe avec $N$ retrouve les deux occupations
par l'équation~\eqref{v:ant:eq:recover}.

La nécessité de la lecture conjointe apparaît dans deux exemples complémentaires.
Les états d'occupation $(1,0)$ et $(0,1)$ ont le même total $N=1$
et un poids modulaire différent. Inversement, $(4,0)$ et $(0,3)$ ont
le même $D=360$ et des totaux différents. La paire $(N,D)$ sépare les deux
cas. Cette récupération concerne les occupations sectorielles, et non
chaque vecteur microscopique au sein d'une dégénérescence de ces occupations.

Lorsque $\Delta_{\rm mod}=0$, l'état est proportionnel à l'identité
et $K_A$ est scalaire : cette lecture cesse de distinguer les secteurs.
Pour $\Delta_{\rm mod}\ne0$, la normalisation connue permet de retrouver
$D$ et de le composer avec l'aire. Le rôle du Hamiltonien
modulaire est ainsi défini : il caractérise l'état réduit et sa distribution d'information.
Son inclusion dans l'article est motivée par cette fonction mathématique,
tandis que le Hamiltonien d'évolution appartient au développement dynamique
de l'action.
